\chapter{Introduzione}
%\chapter*{Introduzione}
%\addcontentsline{toc}{chapter}{Introduzione} % per farla apparire nell’indice
%\markboth{INTRODUZIONE}{}                    % per impostare l'intestazione delle pagine



% 0.1 %%%%%%%%%%%%%%%%%%%%%%%%%%%%%%%%%%%%%%%%%%%%%%%%%%%%%%%%%%%%%%%%%%%%%%%%
\section{Contesto}
%\section*{Contesto}
%\addcontentsline{toc}{section}{Contesto} % per farla apparire nell’indice
L'ultimo decennio, e in particolare il periodo post-pandemico, ha conosciuto una fortissima trasformazione
digitale in ogni settore. Paradigmi come l'Industria 4.0 enfatizzano l'integrazione di tecnologie digitali e 
l’interconnessione nei processi produttivi, creando veri e propri ambienti intelligenti grazie all’adozione di 
strumenti come l’Internet of Things, l’Intelligenza Artificiale e i Big Data.\footcite[pp. 1--2]{clusit-report}
\footcite[pp. 3--4]{enisa-threat-landscape} In questo scenario, le imprese italiane ed europee devono fronteggiare 
nuove minacce alla cybersicurezza e alla protezione delle informazioni,
evolutesi parallelamente al progresso tecnologico.

La sicurezza (in tutte le sue declinazioni) assume dunque un ruolo sempre più centrale nel management aziendale
e deve essere affrontata sempre più attraverso strategie integrate.
Gli attacchi moderni non sono più soltanto fisici o informatici, ma una combinazione di entrambi,
incrementando la complessità della gestione. La valutazione e mitigazione del rischio, la definizione di una soglia 
di rischio accettabile per gli stakeholder e la resilienza delle infrastrutture aziendali rappresentano oggi temi 
cruciali per figure come il CISO (Chief Information Security Officer),
sia a livello di singola organizzazione sia lungo l’intera catena di fornitura.



% 0.2 %%%%%%%%%%%%%%%%%%%%%%%%%%%%%%%%%%%%%%%%%%%%%%%%%%%%%%%%%%%%%%%%%%%%%%%%
\section{Minacce alla sicurezza informatica}
%\section*{Minacce alla sicurezza informatica}
%\addcontentsline{toc}{section}{Minacce alla sicurezza informatica} % per farla apparire nell’indice
Per avere una visione completa sui rischi che le imprese devono gestire e sugli attacchi da fronteggiare,
è utile considerare i principali report di settore, che offrono un quadro aggiornato sulle minacce e sulle
opportunità di rafforzamento dei sistemi di difesa. Tra i riferimenti più autorevoli si annoverano:
\begin{itemize}
    \item \textbf{ENISA Threat Landscape 2024}: rapporto annuale dell'Agenzia dell'Unione europea per la 
        cybersicurezza, che fornisce una panoramica sulle principali minacce, tendenze e vulnerabilità
        a livello europeo.
    \item \textbf{Rapporto Clusit 2025}: pubblicazione dell’Associazione Italiana per la Sicurezza 
        Informatica, che analizza l’andamento degli attacchi informatici e delle tecniche utilizzate
        con focus specifico sul contesto italiano.
\end{itemize}

Entrambi i report rilevano, come si può immaginare, un trend di continua crescita degli episodi di
incidenti informatici rilevati negli ultimi anni.
Inoltre, sottolineano come il quadro geopolitico internazionale contribuisca in modo significativo 
al proliferare di tali eventi. 
Sebbene infatti il \textit{cybercrime} sia costantemente la categoria di attaccanti più prolifica\footcite[p. 22]{clusit-report},
il moltiplicarsi dei conflitti globali e di fenomeni appartenenti alla sfera
dell’\textit{hacktivism} costituiscono fattori sempre più rilevanti nelle statistiche sugli attacchi. 
Tuttavia, la presente tesi si concentrerà esclusivamente sugli aspetti di cybersicurezza in ambito aziendale, 
lasciando sullo sfondo le dimensioni prettamente militari e geopolitiche.

\subsection{Europa}
%\subsection*{Europa}
%\addcontentsline{toc}{subsection}{Europa} % per farla apparire nell’indice
Il \textit{Threat Landscape 2024} di ENISA individua le principali minacce che colpiscono l’ecosistema digitale
europeo. Tra luglio 2023 e giugno 2024 ENISA ha registrato 11,079 incidenti, di cui 322 con un impatto su
più paesi dell'Unione\footcite[p. 12]{enisa-threat-landscape}. DDoS e ransomware si confermano le minacce più
significative dal punto di vista numerico, mentre phishing e social engineering, sempre più potenziati dagli strumenti di intelligenza artificiale, 
risultano in forte crescita sia per volume sia per efficacia (Fig. \ref{fig:incidenti_per_tipo}). 

\defaultimage{enisa-fig2.png}{Ripartizione degli incidenti analizzati per tipologia di minaccia\footcite[Fig. 2]{enisa-threat-landscape}}{fig:incidenti_per_tipo}

A livello europeo, gli attacchi risultano distribuiti trasversalmente tra i diversi settori, 
evidenziando come gli attori malevoli non concentrino le proprie attività su un unico ambito 
specifico (Fig. \ref{fig:incidenti_per_settore}). La cybersicurezza presenta pertanto come 
una sfida che interessa l’intero tessuto economico e sociale, coinvolgendo tutti senza distinzioni rilevanti.

\defaultimage{enisa-fig6.png}{Settori colpiti per numero di incidenti\footcite[Fig. 6]{enisa-threat-landscape}}{fig:incidenti_per_settore}

ENISA evidenzia un incremento degli attacchi alle catene di fornitura, sfruttando i rapporti di fiducia 
tra imprese, fornitori e partner tecnologici come vettori per la propagazione delle minacce attraverso 
le reti.\footcite[pp. 12--13]{enisa-threat-landscape} Attacchi alla supply chain rappresentano una minaccia trasversale 
a più domini, che richiede un approccio ampio alla gestione del rischio, volto al monitoraggio di tutte le dipendenze critiche.

Il report segnala anche la crescita delle compromissioni degli ambienti cloud, sempre più 
utilizzati dalle organizzazioni per la gestione dei dati, a riprova del fatto che gli attacchi data-related restano 
tra i principali obiettivi degli attaccanti, rappresentando una fonte remunerativa di grande rilevanza\footcite[p. 17]{enisa-threat-landscape}.

\subsection{Italia}
%\subsection*{Italia}
%\addcontentsline{toc}{subsection}{Italia} % per farla apparire nell’indice
Il quadro delineato dal \textit{Rapporto Clusit 2025} conferma in larga parte, a livello nazionale,
le tendenze individuate da ENISA. Nel 2024 in Italia si è registrato un aumento del 15,2\% degli 
attacchi rispetto all’anno precedente, mentre la quota di incidenti registrati nel Paese è superiore
del 10\% rispetto al dato globale, segno di una persistente fragilità nella sicurezza dei sistemi 
informatici italiani\footcite[pp. 30--32]{clusit-report}. Un elemento di riflessione in questo senso
è la sproporzione tra il peso economico dell’Italia e l’incidenza degli attacchi: a fronte di un PIL pari a 
circa l’1,8\% di quello mondiale, il Paese concentra il 10,1\% degli incidenti informatici globali\footcite[p. 30]{clusit-report}.
Si segnala comunque che il trend di crescita del numero di incidenti in Italia nel 2024 rimane 
inferiore rispetto al resto del mondo (Fig. \ref{fig:italia_vs_global}). 

\defaultimage{clusit-fig24}{Crescita percentuale in Italia a confronto con il dato globale\footcite[Fig. 24]{clusit-report}}{fig:italia_vs_global}

Un dato significativo riguarda la gravità degli incidenti: gli incidenti classificati come critici
rappresentano solo il 9\% in Italia, contro il dato globale del 29\%. Come sottolineato dal Clusit,
questo dato, da un lato, rappresenta sicuramente un segnale positivo per le infrastrutture IT, ma dall'altro
potrebbe indicare che le organizzazioni italiane sono più spesso vittime di attacchi meno gravi e dunque
presentano una vulnerabilità superiore alla media\footcite[pp. 41--42]{clusit-report}

I dati dell'Osservatorio Cybersecurity del Politecnico di Milano confermano queste ipotesi, evidenziando
come gli incidenti riportati dalle grandi aziende italiane nel 37\% dei casi non abbiano avuto conseguenze
significative\footcite[p. 42]{clusit-report}.



% 0.3 %%%%%%%%%%%%%%%%%%%%%%%%%%%%%%%%%%%%%%%%%%%%%%%%%%%%%%%%%%%%%%%%%%%%%%%%
\section{Standard internazionali e direttive europee}
%\section*{Standard internazionali e direttive europee}
%\addcontentsline{toc}{section}{Standard internazionali e direttive europee} % per farla apparire nell’indice
Per supportare la gestione aziendale e fronteggiare le crescenti minacce informatiche, negli anni sono stati 
sviluppati diversi standard e framework internazionali, con l’obiettivo di promuovere buone pratiche condivise.
Tra i più diffusi si segnalano lo \textit{Standard ISO/IEC 27001:2022} e il \textit{NIST Cybersecurity Framework (CSF) 2.0}, 
oggi strumenti di riferimento per la gestione della sicurezza delle informazioni e del rischio cyber, oltre che per 
l’attuazione di un efficace sistema difensivo\footcite[p. 4]{uni-pdr174}.

La Prassi di Riferimento UNI/PdR 174:2025 riassume efficacemente scopi e differenze dei due standard:
\begin{citazione}[p. 4]{uni-pdr174}
    ISO/IEC 27001 definisce i requisiti per stabilire, attuare, mantenere e migliorare in modo
    continuo un SGSI (Sistema di Gestione per la Sicurezza delle Informazioni), consentendone la
    certificazione volontaria. Tale sistema è basato sulla gestione del rischio e sulla
    determinazione dei controlli necessari a trattarlo.
\end{citazione}
\begin{citazione}[p. 4]{uni-pdr174}
    Il NIST Cybersecurity Framework (CSF) 2.0 è una guida utilizzata da organizzazioni di diversa natura
    per gestire in modo strutturato i rischi di cybersecurity. La guida è centrata su una tassonomia di
    obiettivi che consentono di comprendere, valutare, pianificare, prioritizzare e
    comunicare le attività volte a ridurre l’esposizione ai rischi informatici.
\end{citazione}

Parallelamente, in ambito europeo prosegue il processo normativo sulla cybersicurezza.
La \textit{Direttiva (UE) 2022/2555}, nota come NIS2, è stata adottata dal Parlamento europeo e dal Consiglio
il 14 dicembre 2022, abrogando la precedente \textit{Direttiva (UE) 2016/1148} (NIS). NIS2 introduce misure 
vincolanti per garantire un livello comune elevato di cybersicurezza nell’Unione\footcite[art. 1]{nis2-directive}, 
rappresentando un punto di svolta per le organizzazioni chiamate ad adeguarsi.

Essendo il 2024 l'anno di effettiva entrata in vigore della direttiva e visto il grande numero di
soggetti interessati, soprattutto in Italia, è ancora presto per trarre conclusioni definitive
sul suo impatto nel settore della cybersicurezza. Tuttavia, è ragionevole attendersi che obblighi e sanzioni previsti spingano
le aziende a porre maggiore attenzione alla sicurezza, come è già avvenuto nel settore finanziario, 
uno tra i più normati recentemente con misure quali la DORA, in cui si osserva una riduzione del numero di incidenti
rispetto ai dati storici\footcite[pp. 5--6]{clusit-report}.



% 0.4 %%%%%%%%%%%%%%%%%%%%%%%%%%%%%%%%%%%%%%%%%%%%%%%%%%%%%%%%%%%%%%%%%%%%%%%%
\section{Cosa cambia per le aziende}
%\section*{Cosa cambia per le aziende}
%\addcontentsline{toc}{section}{Cosa cambia per le aziende} % per farla apparire nell’indice
Le implicazioni organizzative dello scenario presentato sono evidenti: l’aumento della superficie d’attacco e la crescente
sofisticazione delle minacce richiedono un approccio integrato di governance e di gestione del rischio. Come
sottolinea il Clusit,
\begin{citazione}[p. 43]{clusit-report}
    La capacità di determinare, anticipare e gestire le evoluzioni legate alle minacce esogene, oltre che al
    contesto interno dell'organizzazione, è ormai fondamentale.
\end{citazione}

La priorità deve essere data alla governance della sicurezza e al presidio dei rischi specifici dell'azienda, 
realizzando misure preventive, mitigative o trasferendo queste responsabilità laddove necessario.
Anche il tema della formazione del personale, soprattutto nelle unità organizzative
esposte a maggior rischio, ricopre un ruolo centrale: l'errore umano rappresenta ancora oggi uno dei 
primi punti di accesso per i cybercriminali alle imprese (basti pensare agli attacchi di phishing,
anch'essi in crescita\footcite[p. 44]{clusit-report}).
La resilienza digitale deve rappresentare un elemento essenziale del “patto di fiducia” tra cittadini, imprese e istituzioni,
obiettivo che costituisce uno dei pilastri fondamentali su cui si erge NIS2\footcite[p. 45]{clusit-report}.



% 0.5 %%%%%%%%%%%%%%%%%%%%%%%%%%%%%%%%%%%%%%%%%%%%%%%%%%%%%%%%%%%%%%%%%%%%%%%%
\section{Obiettivo}
%\section*{Obiettivo}
%\addcontentsline{toc}{section}{Obiettivo} % per farla apparire nell’indice
ISO/IEC 27001, NIST CSF e la Direttiva NIS2 costituiscono i principali riferimenti per la gestione della cybersicurezza aziendale e, 
al contempo, pongono questioni rilevanti di \textit{governance} e responsabilità. Gli standard internazionali forniscono un linguaggio 
comune e un metodo strutturato per organizzare controlli e processi; la Direttiva NIS2 introduce invece obblighi giuridici vincolanti, 
responsabilità dirette del vertice aziendale e un regime sanzionatorio significativo. 

Pur condividendo l’obiettivo di innalzare il livello di sicurezza e resilienza, ciascuno di questi strumenti
presenta approcci e requisiti differenti: gli standard costituiscono una solida base metodologica, ma non
sono di per sé sufficienti a garantire la piena conformità normativa, che richiede l’adozione di misure
aggiuntive e specifiche per la singola azienda.

A partire da tali presupposti, l’elaborato persegue due obiettivi principali:
\begin{enumerate}
  \item Chiarire il rapporto tra strumenti volontari e obblighi normativi: evidenziare in che misura l’adozione di ISO/IEC 27001:2022 e del 
    NIST CSF 2.0 possa agevolare (ma non sostituire) la piena conformità a NIS2, distinguendo le aree di sovrapposizione da quelle che 
    richiedono misure aggiuntive specifiche per la singola impresa.
  \item Analizzare le implicazioni organizzative della direttiva e il ruolo del CISO: indagare come gli obblighi di NIS2 si traducano in responsabilità 
    operative e decisionali per il CISO e per il top management, anche alla luce delle evidenze emerse dalle interviste a professionisti del settore, 
    e come gli standard internazionali possano fungere da strumenti di supporto all’attuazione.
\end{enumerate}

In generale, l'elaborato si propone di analizzare, nella consapevolezza che il tema resti ampio e con l'auspicio
di offrire un contributo utile alla ricerca in materia, come le imprese si adeguino alla Direttiva NIS2, quali nuove
responsabilità e sfide essa comporti per il management e in che modo si rapporti agli standard internazionali menzionati.
La comprensione del ruolo di ISO/IEC 27001 e NIST CSF e delle responsabilità della figura del CISO nel processo di attuazione dei requisiti della 
Direttiva NIS2 costituirà dunque il nucleo centrale dell’analisi sviluppata nei capitoli successivi.



% 0.6 %%%%%%%%%%%%%%%%%%%%%%%%%%%%%%%%%%%%%%%%%%%%%%%%%%%%%%%%%%%%%%%%%%%%%%%%
\section{Metodologia}
%\section*{Metodologia}
%\addcontentsline{toc}{section}{Metodologia} % per farla apparire nell’indice
La metodologia adottata nell’elaborato si articola su tre aspetti, con l’obiettivo di integrare l’analisi 
teorica e documentale con l’evidenza empirica e di restituire una visione completa delle implicazioni organizzative della Direttiva NIS2.

In primo luogo, viene svolta un’analisi della Direttiva NIS2, dello standard \textit{ISO/IEC 27001:2022} e del \textit{NIST Cybersecurity Framework (CSF) 2.0}. 
Tale fase si basa su una revisione sistematica della letteratura scientifica in merito, dei testi normativi e degli standard di riferimento, nonché su 
studi e linee guida elaborate da enti nazionali ed europei (ad esempio ENISA ed ECSO). Lo scopo è evidenziare punti di convergenza e divergenza tra strumenti 
di regolazione vincolante e framework di cybersicurezza, valutando in che misura questi ultimi possano agevolare, ma non sostituire, come si vedrà,
la piena compliance alla normativa europea.

Successivamente, l’elaborato propone un confronto strutturato tra i tre strumenti, con un approccio di \textit{gap analysis}.
Questa parte ha l’obiettivo di rendere immediatamente visibile il grado di copertura dei requisiti della Direttiva da parte degli standard, 
mettendo in luce le aree in cui un’organizzazione deve implementare misure aggiuntive. A tal fine, nel Capitolo \ref{cap:mapping} è presentata una tabella di mappatura che 
delinea l'allineamento delle misure prescritte dall'Art. 21 della Direttiva NIS2 con gli \textit{outcome} esposti nell'allegato del framework NIST
e i controlli dell'\textit{Annex A} di ISO 27001.

Infine, l’analisi condotta viene integrata da un’indagine empirica basata su interviste a Chief Information Security Officer (CISO) di grandi aziende italiane.
Tale approccio consente di cogliere la prospettiva gestionale e strategica delle figure direttamente coinvolte nell’attuazione delle misure di cybersicurezza. 
In particolare, sono stati intervistati:
\begin{itemize}
    \item Matteo Herin, CISO \& Head of AI Cybersecurity presso Carrefour Italia (intervista individuale condotta dall'autore in data 10/09/2025);
    \item Michele Rivieri, CISO di Prometeia (intervista individuale condotta dall'autore in data 17/09/2025).
\end{itemize}
Le interviste hanno l’obiettivo di approfondire il ruolo del CISO, le responsabilità operative e decisionali derivanti dall’attuazione della Direttiva NIS2, 
nonché le modalità con cui gli standard internazionali possono concretamente supportarne l’implementazione in azienda. Le considerazioni emerse vengono 
discusse in relazione alle evidenze tratte dai documenti analizzati, evidenziando criticità, opportunità e \textit{trade-off} organizzativi per il 
top management.

Si segnala che il campione di intervistati, limitato a grandi organizzazioni, non permette di rappresentare pienamente la varietà di realtà 
imprenditoriali interessate dalla Direttiva NIS2, in particolare le PMI, essenza del tessuto imprenditoriale italiano nonché le entità maggiormente 
esposte a vincoli di risorse, fortemente rilevanti in fase di implementazione delle misure di cybersicurezza. Tuttavia, l’obiettivo non è fornire 
un quadro esaustivo di tutte le casistiche, anché perchè arduo, bensì illustrare le principali implicazioni organizzative e gestionali della Direttiva 
attraverso evidenze empiriche e presentare una riflessione critica sul ruolo degli standard internazionali come strumenti di supporto al management.
